\section{Определение суммы и прямой суммы подпространств. Формулировка теоремы о прямой сумме подпространств}

\begin{definition}[Сумма подпространств]
Пусть $U_1, U_2 \subset V$. \textbf{Суммой} подпространств называется
\[
U_1 + U_2 = \{ x \in V : x = x_1 + x_2,\ x_1 \in U_1,\ x_2 \in U_2 \}.
\]

Пусть $U_1 = L(a_1, \ldots, a_k)$, где $E_1 = \{a_1, \ldots, a_k\}$ --- базис в $U_1$.

Пусть $U_2 = L(b_1, \ldots, b_m)$, где $E_2 = \{b_1, \ldots, b_m\}$ --- базис в $U_2$.

Если $a_{i1}, \ldots, a_{ir}, b_{j1}, \ldots, b_{js}$ --- максимальная ЛНЗ система в $E_1 \cup E_2$, тогда $\{a_{i1}, \ldots, a_{ir}, b_{j1}, \ldots, b_{js}\}$ --- базис $U_1 + U_2$.
\end{definition}

\begin{definition}[Прямая сумма]
$U_1 \oplus U_2$ --- \textbf{прямая сумма}, если $\forall x \in U_1 \oplus U_2$ представление $x = x_1 + x_2$ единственно.
\end{definition}

\begin{theorem}[Критерий прямой суммы]
\[
V = U_1 \oplus U_2 \iff U_1 \cap U_2 = \{0\} \quad \text{и} \quad \dim U_1 + \dim U_2 = \dim V.
\]
\end{theorem}

\begin{corollary}[Формула Грассмана]
\[
\dim(U_1 + U_2) = \dim U_1 + \dim U_2 - \dim(U_1 \cap U_2).
\]
\end{corollary}
